\section{Proofs up to reordering}\label{sec:upto}

\paragraph{The problem.} Proc's composition (\file{examples/Proc.v},
\code{compLTS}) writes its structural congruence as silent rules:
\code{do_comm} ($l \mid r \to r \mid l$), \code{do_assocl},
\code{do_assocr} (\cref{sec:containers:upto}). Four processes, each with
three local states, reach $3^4 = 81$ combinations of local states, each
in $120$ arrangements: 5 bracketings of the four components times their
24 orders. That is 9,720 states (\file{examples/Bisimilarity/Proc/Test4}),
and a certificate over states pays 120 times for each combination: over
2.3 million facts, beyond what the proof could hold. Every state of one
combination reaches every other silently, so the 120 are weakly
bisimilar: one of them should stand for all.

\paragraph{Terms.} Write $s \equiv s'$ when $s$ reaches $s'$ by
reorderings; it is an equivalence (each reordering has an inverse). Its
equivalence classes are the \emph{classes}; one state of each, chosen by
the engine, is its \emph{representative}. A shape (\cref{sec:containers})
is the part of a state made of the structural constructors (here
\code{cpar} and \code{cprc}), its \emph{slots} the places where values sit
(here the four processes' terms). A reordering never inspects a value: it
moves the values from slot to slot and may change the shape. So an
\emph{item} of a shape is a shape and a \emph{slot map} (slot $j$ of the
item's shape holds the value of the original's slot $\mathit{map}(j)$),
and its \emph{path} is the reorderings that lead to it.

\paragraph{Why representatives alone are not enough.} Take a state $b$
with one step, silent, to $a$; and $a$ with a silent step back to $b$ and
an $x$-step to $0$. $a$ and $b$ reach each other silently, but $b$ cannot
do $x$ without first moving. A game played on $b$ alone would let an
inert state match the class $\{a, b\}$, which is wrong. What is needed is
that $\equiv$ be a \emph{strong bisimulation}: if $s \equiv s'$ and $s$
steps to $t$, then $s'$ steps, by one step with the same label, to some
$t' \equiv t$. Then weak bisimulation \emph{up to} $\equiv$ is sound
\citep{milner89,pous-sangiorgi11,sangiorgi11}: a game played on
representatives, where every answer may land on any state $\equiv$ to one
in the relation, proves the goal for every member of every class. In the
example the requirement fails ($a \to 0$ by $x$, $b$ has no $x$-step), so
the example is refused.

\paragraph{The engine} (\file{lib/learn/learn.ml}, ``Up to reordering'').
\begin{enumerate}
\item A \emph{reordering summand} is a summand (\cref{sec:containers})
  whose innermost rule is a reordering; its target only moves slots.
\item \emph{Saturation}: from a shape, every reordering summand is fired
  at every place where it applies, to a fixpoint (a state met again stops
  the search, so no orientation of the rules is needed, and
  $l \mid r \to r \mid l$, its own inverse, is no problem). This is done
  once per shape, with no values: the result is the shape's items, each
  with its path. Test4's four-leaf shapes form one class of 120 items.
  A bound (\code{MEBI_UPTO_BOUND}, 50,000 items) stops it.
\item A vector's \emph{representative} is the least, in lexicographic
  order, of its shape's items applied to its values; every member of a
  class chooses the same one.
\item Exploration runs from representatives: every new state is replaced
  by its representative. The game runs on pairs of representatives. A
  move's target is not replaced (the mover's own relations take it
  there), so an answer may bring a state back to its representative at
  any point: \emph{normalisation steps}, one relation per side, are among
  the answer's silent steps. The relations are learned from
  representatives only, which is why an answer goes back to
  representatives as it goes.
\item It applies to weak goals of an LTS whose rules hold reorderings
  and whose summands are not dynamic (a computed value never has
  structure). It is on by default; \code{MEBI_UPTO=off} turns it off.
\end{enumerate}

\emph{Example.} Two processes, \code{cpar (cprc a1) (cprc a2)} against
\code{cpar (cprc a2) (cprc a1)}. The shape \code{cpar (cprc P0) (cprc
P1)} has one reordering summand, \code{do_comm} at the top, so two items:
itself, and itself with the two slots swapped. Each side has 18 states
(9 combinations of local states, each in both orders); up to reordering,
9. The second system's start is not its class's representative: a
normalisation step takes it there.

\paragraph{The theorem} (\file{theories/DdUpto.v}, proved once, closed
under the global context). \code{cert_weak_sim_upto} and
\code{cert_weak_bisim_upto} conclude \code{weak_sim} or
\code{weak_bisimilar} from:
\begin{itemize}
\item each side's equivalence \code{Eq} an equivalence and a strong
  simulation (\code{Equiv});
\item the game on vectors, \code{VWSimN}: every move from a pair of $z$
  answered into $z$, the answer's silent steps the answerer's or
  normalisation steps;
\item the mover's completeness and the answerer's soundness, as in
  \cref{sec:qed};
\item each normalisation step sound: a state, and one \code{Eq} to it
  (\code{SoundNL}, \code{SoundNR});
\item the start normalised into $z$ (\code{NormZ}).
\end{itemize}
Inside, a real answer is carried along \code{Eq} at each normalisation:
if $x \equiv x_0$ and $x_0$ answers by a silent step, so does $x$, to a
state $\equiv$ to $x_0$'s (\code{taus_left_upto}).

\paragraph{The proof} (\file{src/lts_proof.ml}, ``Up to reordering'').
\begin{itemize}
\item \code{Ro}, the one-step reorderings: a disjunction with one case
  per reordering summand of each shape met, over the slots' numbers. A
  slot holds a value of its type's table, \code{T_v u}
  (\cref{sec:qed:parts}). For the example, one case: \code{Ro x y} if
  for some numbers $u_0, u_1$, $x$ is \code{cpar (cprc (T_v $u_0$)) (cprc
  (T_v $u_1$))} and $y$ is the same with $u_0$ and $u_1$ swapped. The
  equivalence is \code{Cong Ro}, its reflexive-transitive closure.
\item Inverse paths: each case's inverse, one step back.
\item \emph{Commutation}, one lemma per case, by the small-inversion
  engine of \cref{sec:qed:parts}: every step of the case's source is one
  of the source's summands. A slot's own step keeps its premise as a
  hypothesis; the reordered state takes the same step at the slot's new
  place, and the targets are again one \code{Ro} step apart; a closure
  lemma (\code{Vn}: each value of the table steps only to values of the
  table) keeps the new value a table value. A reordering of the source is
  answered by the inverse reordering, back to the source. A rule that
  needs a slot to start with a constructor no value of its table starts
  with is refuted by the existing lemma for that (Test4's
  \code{do_par_end}, which needs \code{cprc tend}: no process of Test4 ever
  ends). Then \code{cong_equiv} gives \code{Equiv}.
\item Each item a lemma: the state and the item's are \code{Cong Ro},
  for all numbers, by the item's path.
\item Each normalisation tuple sound: the answerer's soundness of
  \cref{sec:qed:parts} with \code{fun x _ y => Cong Ro x y} in the place
  of the LTS, each tuple's proof its item's lemma at its numbers.
\item The start: one normalisation step per side whose start is not its
  representative (\code{normz_l}, \code{normz_r}).
\end{itemize}
Rocq computes nothing new: the reorderings' targets and the shapes'
templates by conversion, as for the other parts, and the closure and
items by their lemmas' terms. If the proof up to reordering cannot be
built (a rule that does not commute with a reordering, or one with two
premises, which the commutation proofs do not follow), the tactic
decides and proves the goal again without it, with a notice
(\code{Learn.upto_off}).

\paragraph{Cost} (13 November 2026).
\begin{center}
\begin{tabular}{lrrr}
\toprule
goal & states & up to & proof \\
\midrule
Proc, two processes & 18 & 9 & 1.3\,s, 0.5\,GB \\
Proc, three processes & 324 & 27 & 6.2\,s, 1.1\,GB \\
Test4, one goal, tactic & 9,720 & 81 & 52\,s, 3.9\,GB \\
Test4, three goals, \code{MeBi Prove} & 9,720 & 81 & 493\,s, 3.3\,GB \\
\bottomrule
\end{tabular}
\end{center}
Proc Test3's 14 goals are proved up to reordering in the default build.
Test4's 593K facts are mostly the images of the answer's closure
through the normalisation relations, which read and write every
position: narrowing them to the positions a shape's slots use would cut
them. For classes past the bound, sorting (neighbouring values compared
and swapped, each swap an item) is designed but not built; larger systems
are meant to come from compositional proofs (\cref{sec:optimisations}).
